\chapter{Tecnologias}
\label{chap:tecnologias}

\section{Criterio de inventario}
\label{sec:tecnologias-criterio}

Este capitulo documenta las tecnologias confirmadas en el proyecto Robot Explota Globos. La fuente principal es \archivo{requirements.txt}, complementada con \archivo{config/settings.py}, la estructura de aplicaciones y la auditoria tecnica validada. No se incluyen tecnologias no presentes en el repositorio ni componentes de produccion que no hayan sido confirmados.

La plataforma resuelve tres necesidades tecnicas principales: recibir informacion publica de participantes, administrar internamente una competencia con reglas variables y emitir comunicaciones con documentos personalizados. Las tecnologias seleccionadas en el proyecto se agrupan alrededor de esas necesidades: Django para estructura web y persistencia, dependencias Python especializadas para documentos y Excel, assets estaticos para interfaz, y servicios externos configurables para correo y conversion PDF.

\section{Stack confirmado}
\label{sec:tecnologias-stack}

\begin{longtable}{p{3.3cm}p{4cm}p{6.2cm}}
\toprule
Capa & Tecnologia confirmada & Proposito dentro de la plataforma \\
\midrule
Backend & Python con Django & Ejecuta vistas, formularios, servicios, autenticacion, administracion, ORM, comandos de gestion y renderizado de templates. \\
ORM & Django ORM & Persiste registros, equipos, competencias, grupos, batallas, estados, historial, plantillas, destinatarios, lotes y logs. \\
Base local & SQLite en desarrollo & Permite preparar y validar el proyecto localmente cuando \variable{DJANGO_DATABASE=sqlite} y \variable{DJANGO_DEBUG=True}. \\
Base no local & PostgreSQL & Configuracion prevista para entornos no locales mediante variables \variable{POSTGRES_*}. \\
Frontend & Django Templates, HTML, CSS y JavaScript & Renderiza vistas publicas e internas sin introducir una aplicacion cliente separada. \\
Assets & Directorio \archivo{static/} & Contiene CSS, JavaScript, PDF publico del reglamento y logos de patrocinadores. \\
Correo & \archivo{django.core.mail}, SMTP o consola & Soporta comunicaciones, pruebas controladas y envios reales segun variables de entorno. \\
Documentos & DOCX y PDF & Usa plantillas DOCX para comunicaciones y convierte salidas a PDF mediante LibreOffice. \\
Excel & XLSX & Carga destinatarios para comunicaciones desde archivos de hoja de calculo. \\
\bottomrule
\end{longtable}

La interaccion entre estas capas aparece en los flujos centrales del sistema. Un registro publico entra por una vista Django, se valida mediante formularios, se persiste con el ORM y despues puede activar servicios que crean equipos y competencias. Un envio de invitaciones parte de plantillas DOCX y destinatarios XLSX, genera documentos personalizados, invoca LibreOffice para PDF y registra resultados en modelos de comunicaciones. El panel interno usa templates y JavaScript para enviar acciones AJAX a vistas Django, que delegan cambios en servicios de torneo.

\section{Dependencias Python}
\label{sec:tecnologias-dependencias-python}

Las dependencias declaradas en \archivo{requirements.txt} son reducidas y estan alineadas con funcionalidades verificadas en el codigo. Esta seleccion favorece mantenimiento porque evita un ecosistema amplio de paquetes para tareas que Django ya cubre.

\begin{longtable}{p{3.2cm}p{3cm}p{7.2cm}}
\toprule
Dependencia & Version declarada & Uso confirmado \\
\midrule
\archivo{Django} & \variable{>=5.2,<5.3} & Framework principal del sistema: URLs, vistas, templates, formularios, ORM, autenticacion, admin, seguridad base y comandos de gestion. \\
\archivo{psycopg[binary]} & \variable{>=3.2,<3.3} & Driver PostgreSQL usado cuando la configuracion no selecciona SQLite local. \\
\archivo{python-dotenv} & \variable{>=1.0,<2.0} & Carga el archivo \archivo{.env} desde \archivo{config/settings.py} para separar configuracion y secretos del codigo. \\
\archivo{openpyxl} & \variable{>=3.1,<3.2} & Permite leer archivos XLSX de destinatarios en el modulo de comunicaciones. \\
\archivo{python-docx} & \variable{>=1.1,<1.2} & Permite inspeccionar marcadores y renderizar plantillas DOCX usadas para invitaciones y solicitudes. \\
\bottomrule
\end{longtable}

La version exacta de Python validada en el entorno virtual local es 3.12.10. La version exacta de Python en produccion permanece como \MarcadorProduccionPendiente. Por esa razon, los capitulos de despliegue y operacion deben validar compatibilidad del servidor final antes de usar esta documentacion como procedimiento productivo cerrado.

\section{Django dentro del proyecto}
\label{sec:tecnologias-django-proyecto}

Django se usa como plataforma de integracion, no solo como capa HTTP. En este proyecto aporta:

\begin{itemize}
  \item \textbf{Rutas y vistas}: \archivo{config/urls.py} distribuye solicitudes hacia \archivo{apps.core}, \archivo{apps.participants}, \archivo{apps.tournament} e \archivo{apps.invitations}.
  \item \textbf{ORM}: los modelos persistidos representan inscripciones, equipos, competencias, batallas, estados, historial, plantillas y logs.
  \item \textbf{Forms}: las validaciones de registro y comunicaciones se concentran en formularios Django, lo que evita dispersar reglas basicas en templates o JavaScript.
  \item \textbf{Templates}: las interfaces publicas e internas se renderizan en servidor y se complementan con JavaScript puntual.
  \item \textbf{Autenticacion y autorizacion}: las areas internas usan login y validacion de usuario organizador.
  \item \textbf{Admin}: se usa para inspeccion y operaciones administrativas sobre participantes, torneo e invitaciones.
  \item \textbf{Comandos}: existen comandos de gestion para bootstrap de torneo, demo, invitaciones y solicitudes de asistencia.
\end{itemize}

Esta concentracion en Django reduce la cantidad de componentes externos necesarios para operar la plataforma. La contrapartida es que los cambios deben respetar las fronteras internas del proyecto: vistas para entrada HTTP, servicios para regla de negocio, modelos para persistencia y templates/JavaScript para presentacion.

\section{Persistencia}
\label{sec:tecnologias-persistencia}

La configuracion distingue desarrollo local y entornos no locales. SQLite se permite solo cuando el proyecto esta en modo desarrollo. Si \variable{USE_SQLITE=True} y \variable{DEBUG=False}, \archivo{config/settings.py} lanza \clase{ImproperlyConfigured}. Esta restriccion evita que una instancia con configuracion de produccion arranque accidentalmente sobre una base local.

Cuando no se usa SQLite, la configuracion exige \variable{POSTGRES_DB}, \variable{POSTGRES_USER}, \variable{POSTGRES_PASSWORD} y \variable{POSTGRES_HOST}. PostgreSQL esta soportado por \archivo{psycopg[binary]}. El puerto se toma de \variable{POSTGRES_PORT} y el valor de ejemplo es 5432.

\begin{advertencia}
El soporte PostgreSQL esta preparado en el codigo, pero el servidor final de produccion, proveedor, host, politica de backups y parametros reales siguen como \MarcadorProduccionPendiente. No se deben copiar al manual valores de un archivo \archivo{.env} real.
\end{advertencia}

\section{Frontend y assets}
\label{sec:tecnologias-frontend-assets}

El frontend se organiza con templates Django y archivos estaticos. La auditoria confirma:

\begin{itemize}
  \item \archivo{static/css/app.css}: estilos publicos e internos.
  \item \archivo{static/js/app.js}: tema, accesibilidad del boton de tema, comportamiento del header y animaciones publicas.
  \item \archivo{static/js/control.js}: interacciones del panel interno, modales, resultados, guardado AJAX y refresco parcial.
  \item \archivo{static/docs/ExplotaGlobos.pdf}: documento publico embebido.
  \item \archivo{static/sponsors/}: logos de patrocinadores usados por el sitio publico.
\end{itemize}

No se observa un gestor de paquetes frontend ni un bundler dentro del repositorio. Esta decision simplifica la preparacion local: no requiere instalar Node.js ni ejecutar compilaciones de assets para que el proyecto funcione. El costo tecnico es que los contratos entre templates y \archivo{control.js} deben mantenerse manualmente, especialmente los atributos \variable{data-*} usados por el panel.

\section{Comunicaciones y documentos}
\label{sec:tecnologias-comunicaciones}

El modulo de comunicaciones usa tres piezas externas al nucleo Django:

\begin{itemize}
  \item \archivo{python-docx}, para leer plantillas DOCX y reemplazar marcadores.
  \item \archivo{openpyxl}, para importar destinatarios desde XLSX.
  \item LibreOffice en modo headless, para convertir documentos DOCX renderizados a PDF.
\end{itemize}

La conversion PDF no depende de una libreria Python pura. El servicio busca LibreOffice por PATH o por \variable{LIBREOFFICE_BINARY}. Esta dependencia es aceptable para el caso de uso porque el sistema necesita generar adjuntos PDF desde plantillas editables, pero debe validarse antes de usar comunicaciones reales en otro ambiente.

\section{Tecnologias no confirmadas}
\label{sec:tecnologias-no-confirmadas}

\begin{longtable}{p{5cm}p{8.2cm}}
\toprule
Elemento & Estado documental \\
\midrule
Servidor WSGI o ASGI de produccion & \MarcadorProduccionPendiente \\
Proveedor de hosting & \MarcadorProduccionPendiente \\
Version exacta de Python en produccion & \MarcadorProduccionPendiente \\
Servicio SMTP real & Configuracion soportada; valores reales no confirmados \\
Binario LibreOffice en entorno final & Requerido para PDF; ubicacion no confirmada \\
Servicio final de archivos estaticos & \MarcadorProduccionPendiente \\
Servicio final de media & \MarcadorProduccionPendiente \\
\bottomrule
\end{longtable}

\section{Cierre del capitulo}
\label{sec:tecnologias-cierre}

El stack del proyecto es deliberadamente compacto: Django cubre la mayor parte de la aplicacion, y las dependencias adicionales responden a necesidades especificas de PostgreSQL, configuracion por entorno, Excel, DOCX y PDF. Esta composicion favorece la mantenibilidad porque cada paquete tiene una responsabilidad clara. Los principales puntos de atencion no estan en la cantidad de tecnologias, sino en las dependencias operativas externas: PostgreSQL en entornos no locales, SMTP real y LibreOffice.
